\documentclass[tikz,border=6pt]{standalone}
\usepackage{ctex}
\usepackage{amsmath}
\usetikzlibrary{calc, arrows.meta}

\begin{document}
\begin{tikzpicture}[
  scale=1.2, thick,
  vec/.style={-{Stealth[length=7pt, width=4pt]}, thick},
]

% 原点
\coordinate (O) at (0, 0);

% 向量 a：水平向右
\coordinate (A) at (3.2, 0);

% 向量 b：竖直向上（与 a 垂直）
\coordinate (B) at (0, 2.5);

% 直角标记（在原点处）
\draw[black, thin]
  (0.18, 0) -- (0.18, 0.18) -- (0, 0.18);

% 画向量
\draw[vec, black, very thick] (O) -- (A) node[right] {$\vec{a}$};
\draw[vec, red, very thick]   (O) -- (B) node[above] {$\vec{b}$};

% 原点标记
\filldraw[black] (O) circle (1.5pt) node[below left] {$O$};
\filldraw[black] (A) circle (1.2pt);
\filldraw[black] (B) circle (1.2pt);

% 垂直符号标注
\node[font=\small] at (1.8, 1.5)
  {$\vec{a} \perp \vec{b}$};

% 数量积公式说明
\node[font=\small, align=center] at (1.6, -0.65)
  {$\theta = 90°$，$\cos 90° = 0$};
\node[font=\small, align=center] at (1.6, -1.15)
  {$\vec{a}\cdot\vec{b} = |\vec{a}||\vec{b}|\cos 90° = 0$};

\end{tikzpicture}
\end{document}
